\reportelement{Введение}

\textbf{Цель работы:} перенести запуск приложения, PostgreSQL и Redis в
Docker Compose и освоить масштабирование сервисов.

В работе используются следующие инструменты: Docker, Docker Compose,
PostgreSQL, Redis, а также приложение-ассистент, обращающееся к внешнему
API погоды. Персональные параметры приведены в
таблице~\ref{tab:parameters}.

\begin{table}[H]
  \caption{Персональные параметры}\label{tab:parameters}
  \centering
  \begin{tabularx}{\textwidth}{|p{0.46\textwidth}|X|}
    \hline
    Параметр & Значение\\ \hline
    Номер ИСУ & \StudentISU\\ \hline
    Персональный порт (PostgreSQL, \texttt{PG\_PORT}) & 5500\\ \hline
    Персональный порт (Redis, \texttt{REDIS\_PORT}) & 6440\\ \hline
    Персональный порт (приложение, \texttt{APP\_PORT}) & 8070\\ \hline
    Персональный префикс & \texttt{lyasin}\\ \hline
    Секретный ключ & \texttt{secret70lyasin}\\ \hline
    Имя базы данных & \texttt{lyasin70}\\ \hline
  \end{tabularx}
\end{table}

Результат хранится в ветке \texttt{lab-\LabNumber} форка репозитория:
\url{\RepositoryURL}.

\chapter{Ход выполнения}

\section{Этап 1. Подготовка проекта для Docker Compose}

Создадим файл \texttt{.env} с переменными окружения проекта
(рисунок~\ref{fig:stage1-env}). Создадим файл \texttt{compose.yml} в корне
проекта: в нём описаны два сервиса~--- база данных и приложение;
переменные окружения подтягиваются из \texttt{.env}
(рисунок~\ref{fig:stage1-compose}). Далее поднимем его командой \texttt{docker compose up -d}

\reportscreenshot{image-1.png}{Файл \texttt{.env} с переменными окружения}{fig:stage1-env}
\reportscreenshot{image-2.png}{Создание файла \texttt{compose.yml}}{fig:stage1-compose}

Проверим логи приложения и базы данных
(рисунки~\ref{fig:stage1-logs-app} и~\ref{fig:stage1-logs-db}).

\reportscreenshot{image-3.png}{Логи приложения}{fig:stage1-logs-app}
\reportscreenshot{image-4.png}{Логи базы данных}{fig:stage1-logs-db}

Убедимся, что сервис успешно поднялся (рисунок~\ref{fig:stage1-up}).

\reportscreenshot{image-5.png}{Проверка успешного запуска сервиса}{fig:stage1-up}

\section{Этап 2. Добавление Redis для кэширования погоды}

Добавим Redis в проект, чтобы кэшировать запросы к внешнему API. Пропишем
новый сервис \texttt{redis} и добавим его в зависимости сервиса
\texttt{app} (рисунок~\ref{fig:stage2-redis-service}).

\reportscreenshot{image-6.png}{Добавление сервиса Redis и зависимостей}{fig:stage2-redis-service}

Файл \texttt{main.py} устарел, поэтому в него были внесены изменения для
работы кэширования. Перезапустим Compose и проверим, что все три сервиса
подняты и имеют статус \texttt{healthy} (кроме \texttt{app}, для него
healthcheck не прописан), рисунок~\ref{fig:stage2-ps}.

\reportscreenshot{image-7.png}{Статусы сервисов после перезапуска}{fig:stage2-ps}

Проверим порядок запуска сервисов через логи: \texttt{app} запускается
последним (рисунок~\ref{fig:stage2-order}).

\reportscreenshot{image-8.png}{Порядок запуска сервисов по логам}{fig:stage2-order}

На этом этапе обнаружено, что подключение к Redis не срабатывает. В
\texttt{compose.yml} добавлена строка, передающая переменную окружения
\texttt{REDIS\_URL} в \texttt{app} (рисунок~\ref{fig:stage2-redis-url}).

\reportscreenshot{image-9.png}{Передача \texttt{REDIS\_URL} в приложение}{fig:stage2-redis-url}

Выведем лог \texttt{app} и убедимся, что подключение к Redis успешно
(рисунок~\ref{fig:stage2-redis-ok}).

\reportscreenshot{image-10.png}{Успешное подключение к Redis}{fig:stage2-redis-ok}

Также \texttt{main.py} был отредактирован для использования зеркал, так
как исходный API оказался недоступен. Выполним два запроса к приложению:
второй произошёл заметно быстрее, в поле \texttt{source} первого запроса
указано \texttt{api}, второго~--- \texttt{cache}
(рисунок~\ref{fig:stage2-cache}).

\reportscreenshot{image-11.png}{Два запроса: источник \texttt{api} и \texttt{cache}}{fig:stage2-cache}

Проверим статистику кэша (рисунок~\ref{fig:stage2-stats}) и Redis
напрямую~--- убедимся в наличии записи
(рисунок~\ref{fig:stage2-redis-direct}).

\reportscreenshot{image-12.png}{Статистика кэша}{fig:stage2-stats}
\reportscreenshot{image-13.png}{Запись в Redis, полученная напрямую}{fig:stage2-redis-direct}

Посмотрим список сетей: созданная сеть имеет приставку
\texttt{assist-py\_} (рисунок~\ref{fig:stage2-networks}). Информация об
этой сети показывает, что все три контейнера принадлежат ей
(рисунок~\ref{fig:stage2-network-info}).

\reportscreenshot{image-14.png}{Список сетей}{fig:stage2-networks}
\reportscreenshot{image-15.png}{Информация о сети проекта}{fig:stage2-network-info}

Аналогично имя volume также имеет приставку \texttt{assist-py\_}
(рисунок~\ref{fig:stage2-volume}).

\reportscreenshot{image-16.png}{Именованный volume проекта}{fig:stage2-volume}

\section{Этап 3. Масштабирование сервиса}

Закоммитим текущий \texttt{compose.yml} и обновим его, добавив возможность
масштабирования (рисунок~\ref{fig:stage3-compose}).

\reportscreenshot{image-17.png}{Обновлённый \texttt{compose.yml} с поддержкой масштабирования}{fig:stage3-compose}

Остановим текущие контейнеры и запустим обновлённый Compose с количеством
экземпляров \texttt{app = 3}. Каждому из трёх контейнеров приложения
назначен свой порт из диапазона 32768--32770
(рисунок~\ref{fig:stage3-scale}).

\reportscreenshot{image-18.png}{Запуск трёх экземпляров приложения}{fig:stage3-scale}

Проверим, что каждый из сервисов доступен по отдельности и что все
экземпляры разделяют общий кэш (рисунок~\ref{fig:stage3-shared}).

\reportscreenshot{image-19.png}{Доступность экземпляров и общий кэш}{fig:stage3-shared}

После проверки вернём предыдущую конфигурацию без масштабирования.

\section{Этап 4. Финальный артефакт}

Финальный артефакт лабораторной работы приведён на
рисунке~\ref{fig:final}.

\reportscreenshot{image-20.png}{Финальный артефакт лабораторной работы}{fig:final}

\reportelement{Заключение}

В ходе работы запуск приложения, PostgreSQL и Redis были выполнены в
Docker Compose. настроены зависимости и порядок запуска. Добавлен
Redis для кэширования запросов к внешнему API. Проверено совместное
использование сети и именованного volume сервисами проекта, а также
выполнено масштабирование приложения до трёх экземпляров с общим кэшем.